\documentclass[20pt,a4paper]{article}
%-----------------------------------------------%
% หมายเหตุการเขียนภาษาไทย
% - XeTeX ตัดคำไทยเอง (\XeTeXlinebreaklocale "th") ไม่ต้องใส่ \\ หรือ \newline กลางย่อหน้า
% - การขึ้นบรรทัดใหม่ในไฟล์ .tex จะกลายเป็นช่องว่าง 1 ช่อง จึงขึ้นบรรทัดใหม่เฉพาะตรงที่
%   ภาษาไทยเว้นวรรคอยู่แล้ว (ระหว่างวลี) ห้ามขึ้นบรรทัดกลางคำ
% - ใช้ ~ (ช่องว่างที่ห้ามตัดบรรทัด) ระหว่างตัวเลขกับหน่วย และหน้า \ref / \cite
% - ย่อหน้าใหม่ใช้บรรทัดว่าง 1 บรรทัด
% - ตัวเลขผลการทดลองทั้งหมดมาจากภาพสังเคราะห์ (synthetic) ยังไม่ได้ทดสอบกับภาพจริง
%-----------------------------------------------%
%ชื่อหัวข้อ (ภาษาไทย)
\newcommand{\titleTH}{การแบ่งส่วนวัตถุระดับ 4~มุมของรหัส DataMatrix ขนาดเล็กหลายรหัสในภาพเดียว เพื่ออ่านฉลากสไลด์พยาธิวิทยาอัตโนมัติ}
%ชื่อหัวข้อ (ภาษาอังกฤษ)
\newcommand{\titleEN}{Corner-Level Instance Segmentation of Multiple Small DataMatrix Codes for Automated Reading of Pathology Slide Labels}

%รหัสนักศึกษาคนที่ 1
\newcommand{\firstStudentId}{673380353-1}
%คำนำหน้า (ภาษาไทย) นักศึกษาคนที่ 1
\newcommand{\firstStudentPrefixTH}{นาย}
%ชื่อ-สกุล (ภาษาไทย) นักศึกษาคนที่ 1
\newcommand{\firstStudentNameTH}{อชิรวิทย์ ศรีชา}
%คำนำหน้า (ภาษาอังกฤษ) นักศึกษาคนที่ 1
\newcommand{\firstStudentPrefixEN}{Mr.}
%ชื่อ-สกุล (ภาษาอังกฤษ) นักศึกษาคนที่ 1
\newcommand{\firstStudentNameEN}{Achirawit Sricha}

%รหัสนักศึกษาคนที่ 2
\newcommand{\secondStudentId}{673380521-6}
%คำนำหน้า (ภาษาไทย) นักศึกษาคนที่ 2
\newcommand{\secondStudentPrefixTH}{นาย}
%ชื่อ-สกุล (ภาษาไทย) นักศึกษาคนที่ 2
\newcommand{\secondStudentNameTH}{ธีรปรัชญ์ สาขาคำ}
%คำนำหน้า (ภาษาอังกฤษ) นักศึกษาคนที่ 2
\newcommand{\secondStudentPrefixEN}{Mr.}
%ชื่อ-สกุล (ภาษาอังกฤษ) นักศึกษาคนที่ 2
\newcommand{\secondStudentNameEN}{Teeraprat Sakhakham}
\newcommand{\major}{ปัญญาประดิษฐ์}

\newcommand{\advisor}{ผศ.ดร.~ไพรสันต์ ผดุงเวียง}

%-----------------------------------------------%
\usepackage[a4paper, top=2.5cm, bottom=2.5cm, left=2.5cm, right=2.5cm]{geometry}
\usepackage{amsmath,esint,mathtools}
\usepackage{indentfirst}
\usepackage{enumitem}
\usepackage{adjustbox}
\usepackage{float}
\usepackage{hyperref}
\usepackage{xcolor}
\usepackage{graphicx}
\usepackage{caption}
\usepackage{subcaption}
\usepackage{xltxtra}
\usepackage{multirow}
\usepackage{array}
\usepackage{titlesec}   % for \titleformat, \titlespacing
\usepackage{etoolbox}   % for \pretocmd
%------------ สำหรับการทำ Gantt Chart -----------%
\usepackage{colortbl}
\usepackage{forloop}
\newcounter{loopcntr}
\newcommand{\rpt}[2][1]{%
  \forloop{loopcntr}{0}{\value{loopcntr}<#1}{#2}%
}
\newcommand{\on}[1][1]{
  \forloop{loopcntr}{0}{\value{loopcntr}<#1}{&\cellcolor{lightgray}}
}
\newcommand{\off}[1][1]{
  \forloop{loopcntr}{0}{\value{loopcntr}<#1}{&}
}



\newlength{\firstindentlength}
\newlength{\secondindentlength}
\newlength{\thirdindentlength}

\setlength{\firstindentlength}{3em}
\setlength{\secondindentlength}{7.1em}
\setlength{\thirdindentlength}{12.2em}

%------------ สำหรับการใช้ฟ้อนท์ภาษาไทย -----------%
\usepackage[no-math]{fontspec} % สำหรับเลือก font
\usepackage{polyglossia}     % สำหรับหลายภาษา
\usepackage{xunicode}

\XeTeXlinebreaklocale "th"
\XeTeXlinebreakskip = 0pt plus 1pt

\renewcommand{\baselinestretch}{1.6}
\setmainfont[Script=Thai,Scale=1.4,ItalicFont={THSarabunNew_Italic.ttf}, BoldFont={THSarabunNew_Bold.ttf},BoldItalicFont={THSarabunNew_BoldItalic.ttf} ]{THSarabunNew.ttf}
\setmonofont[Script=Thai,Scale=0.9]{DroidSansMono.ttf}
%--------------- เปลี่ยนชื่อรูปและตาราง -------------%
\usepackage[labelsep=space]{caption}
\usepackage{enumitem}

\captionsetup{
  figurename=ภาพที่,
  tablename=ตารางที่
}

%-----------------------------------------------%
\setlength{\parindent}{\firstindentlength}
\setlist[enumerate]{
label=\thesection.\arabic*,
leftmargin=\parindent,
labelsep=1em,
itemindent=\parindent,
}
\makeatother
\setlist[itemize]{
leftmargin=\parindent,
labelsep=1em,
itemindent=\parindent,
}
\setlist[description]{
leftmargin=\parindent,
labelsep=1em,
itemindent=0.5\parindent,
}
%-----------------------------------------------%
% \setlist[itemize]{%
%   label=\textbullet
% }
% \setlist[enumerate]{label*=\arabic*.}

% \setlist[itemize,1]{%
%   leftmargin=*,
%   labelindent=\parindent,
%   labelsep=1em,
%   listparindent=0pt
% }
% \setlist[itemize,2]{%
%   leftmargin=*,
%   labelindent=0em,
%   labelsep=1em,
%   listparindent=0pt
% }
% \setlist[itemize,3]{%
%   leftmargin=*,
%   labelindent=0em,
%   labelsep=1em,
%   listparindent=0pt
% }

% \setlist[enumerate,1]{%
%   label=\thesection.\arabic*.,
%   leftmargin=*,
%   labelindent=\parindent,
%   labelsep=1em,
%   listparindent=0pt
% }
% \setlist[enumerate,2]{%
%   label=\thesection.\arabic{enumi}.\arabic*.,
%   leftmargin=*,
%   labelindent=0em,
%   labelsep=1em,
%   listparindent=0pt
% }
% \setlist[enumerate,3]{%
%   label=\thesection.\arabic{enumi}.\arabic{enumii}.\arabic*.,
%   leftmargin=*,
%   labelindent=0em,
%   labelsep=1em,
%   listparindent=0pt
% }

% \setcounter{secnumdepth}{3}
% \renewcommand\thesection{\arabic{section}}
% \renewcommand\thesubsection{\thesection.\arabic{subsection}}
% \renewcommand\thesubsubsection{\thesubsection.\arabic{subsubsection}}

% \titleformat{\section}
%   {\bfseries\large\upshape}
%   {\thesection}
%   {0.45\firstindentlength}
%   {}

% \titleformat{name=\section,numberless}
%   {\bfseries\large\upshape}
%   {}
%   {0pt}
%   {}

% \titleformat{\subsection}
%   {\bfseries\large\upshape}
%   {\thesubsection}
%   {0.45\firstindentlength}
%   {}

% \titleformat{name=\subsection,numberless}
%   {\bfseries\large\upshape}
%   {}
%   {0pt}
%   {}

% \titleformat{\subsubsection}
%   {\bfseries\large\upshape}
%   {\thesubsubsection}
%   {0.45\firstindentlength}
%   {}

% \titleformat{name=\subsubsection,numberless}
%   {\bfseries\large\upshape}
%   {}
%   {0pt}
%   {}

% \titlespacing{\section}
%   {0pt}
%   {\baselineskip}
%   {0\baselineskip}

% \titlespacing{\subsection}
%   {\firstindentlength}
%   {\baselineskip}
%   {0\baselineskip}

% \titlespacing{\subsubsection}
%   {\secondindentlength}
%   {\baselineskip}
%   {0\baselineskip}

% \pretocmd{\section}{\setlength{\parindent}{\firstindentlength}}{}{}
% \pretocmd{\subsection}{\setlength{\parindent}{\secondindentlength}}{}{}
% \pretocmd{\subsubsection}{\setlength{\parindent}{\thirdindentlength}}{}{}

\begin{document}
\begin{center}
\Large
\textbf{เค้าโครงโครงงานปัญญาประดิษฐ์และรายงานความก้าวหน้า}
\\
\large
วิทยาลัยการคอมพิวเตอร์ มหาวิทยาลัยขอนแก่น
\end{center}

\noindent
\begin{tabular}{@{}p{0.5\textwidth}@{}@{}p{0.5\textwidth}@{}}
\\
\textbf{ชื่อ (ภาษาไทย)} \firstStudentPrefixTH\firstStudentNameTH
& \textbf{รหัสนักศึกษา} \firstStudentId
\\ \textbf{ชื่อ (ภาษาอังกฤษ)} \firstStudentPrefixEN\ \firstStudentNameEN  &
\\
 &
\\
\textbf{ชื่อ (ภาษาไทย)} \secondStudentPrefixTH\secondStudentNameTH
& \textbf{รหัสนักศึกษา} \secondStudentId
\\ \textbf{ชื่อ (ภาษาอังกฤษ)} \secondStudentPrefixEN\ \secondStudentNameEN  &
\\
\end{tabular}

\vspace{1em}
\noindent \textbf{นักศึกษาระดับปริญญาตรี} หลักสูตรวิทยาศาสตรบัณฑิต สาขาวิชา\major


\vspace{1em}
\noindent \textbf{อาจารย์ที่ปรึกษาโครงงาน} \advisor

\noindent\rule{\textwidth}{0.1pt}

%-----------------------------------------------%
% 1. Research Title
%-----------------------------------------------%

\section{ชื่อหัวข้อโครงงาน (Research Title)}

\begin{quote}
\textbf{ชื่อหัวข้อ (ภาษาไทย):} \titleTH
\end{quote}

\begin{quote}
\textbf{ชื่อหัวข้อ (ภาษาอังกฤษ):} \titleEN
\end{quote}

%-----------------------------------------------%
% 2. Introduction and Motivation
%-----------------------------------------------%

\section{หลักการและเหตุผล
(Introduction and Motivation)}

\subsection{ความเป็นมา
(Background and Motivation)}

ในงานพยาธิวิทยา ชิ้นเนื้อของผู้ป่วยแต่ละรายถูกเตรียมเป็นสไลด์แก้วจำนวนมาก
และสไลด์ทุกแผ่นต้องผูกกับเลขที่ตัวอย่างของผู้ป่วยให้ถูกต้องตลอดกระบวนการ
การระบุตัวอย่างผิดจึงเป็นปัญหาด้านความปลอดภัยของผู้ป่วยโดยตรง
Zarbo และคณะ~\cite{zarbo2009henryford} รายงานว่า
เมื่อห้องปฏิบัติการพยาธิวิทยาปรับกระบวนการทำงานร่วมกับการติดฉลากบาร์โค้ดบนสไลด์
อัตราการระบุตัวอย่างผิดโดยรวมลดลงประมาณ~62\%
และความผิดพลาดของสไลด์ทางจุลพยาธิวิทยาลดลงประมาณ~95\%

แม้ฉลากสไลด์จะมีรหัส DataMatrix แล้ว
แต่ในทางปฏิบัติเจ้าหน้าที่ยังต้องสแกนรหัสทีละแผ่นด้วยเครื่องอ่านมือถือ
ซึ่งใช้เวลามากเมื่อมีสไลด์จำนวนมาก
แนวคิดของโครงงานนี้คือถ่ายภาพสไลด์ทั้งแถว (สูงสุด 10~แผ่น) ในภาพเดียว
แล้วใช้ปัญญาประดิษฐ์หาตำแหน่งและอ่านรหัสทุกอันพร้อมกัน
โมเดลตรวจจับวัตถุรุ่นใหม่ เช่น YOLO~\cite{khanam2024yolov11,sapkota2025yolo26}
และ RF-DETR~\cite{robinson2026rfdetr}
ทำงานได้เร็วพอสำหรับการใช้งานแบบ real-time
จึงมีศักยภาพที่จะนำมาใช้กับงานนี้

\subsection{ปัญหาวิจัยและข้อจำกัดทางเทคนิค
(Research Problems and Technical Challenges)}

เมื่อถ่ายสไลด์ 10~แผ่นด้วยกล้องความละเอียด 8--12~ล้านพิกเซล
รหัส DataMatrix ขนาด 18$\times$18~module กว้าง 6.5~มม.
จะมีความละเอียดในภาพเพียงประมาณ 3.5--5~พิกเซลต่อ module (px/module)
และรหัสอาจเอียงหรือบิดตามมุมกล้อง
ทำให้เกิดข้อจำกัดดังนี้

\begin{enumerate}
    \item \textbf{ตัวอ่านสำเร็จรูปอ่านทั้งภาพได้ไม่ครบ:}
    ผลการทดลองเบื้องต้นบนภาพสังเคราะห์พบว่า
    เมื่อส่งภาพทั้งภาพให้ zxing-cpp~\cite{zxingcpp} โดยตรง
    อ่านได้เพียง 64.1\% ที่ช่วงความละเอียดของกล้อง 8~ล้านพิกเซล (3.5--4.0~px/module)

    \item \textbf{กรอบสี่เหลี่ยมไม่แก้ความเอียง:}
    งานที่ตรวจจับแล้วถอดรหัสจากกรอบสี่เหลี่ยม~\cite{alsulami2026uav}
    รายงานว่ารหัสเอียงแล้วอ่านได้ลดลง

    \item \textbf{ยังไม่รู้ความละเอียดขั้นต่ำ:}
    งานที่เกี่ยวข้องไม่ได้รายงานผลแยกตามความละเอียดของรหัส
    จึงยังไม่รู้ว่าภาพต้องละเอียดเท่าใด

    \item \textbf{ไม่มีชุดข้อมูลที่ใช้ได้:}
    ภาพฉลากสไลด์จริงมีเลขที่ตัวอย่างผู้ป่วย
    ซึ่งเป็นข้อมูลสุขภาพตาม พ.ร.บ.~คุ้มครองข้อมูลส่วนบุคคล (PDPA)
    จึงเผยแพร่เป็นชุดข้อมูลไม่ได้
    และการติดป้ายกำกับ 4~มุมด้วยมือในภาพความละเอียดต่ำก็ทำได้ยากและไม่แม่นยำ
\end{enumerate}

%-----------------------------------------------%
% 3. Literature Review
%-----------------------------------------------%

\section{ทฤษฎีและผลงานวิจัยที่เกี่ยวข้อง (Literature Review)}

\subsection{ทฤษฎีที่เกี่ยวข้อง (Theoretical Background)}

\textbf{รหัส DataMatrix และความละเอียดของภาพ}
DataMatrix เป็นรหัส 2~มิติที่ประกอบด้วยช่องสี่เหลี่ยมเล็ก (module)
มีขอบรูปตัว L ใช้หาตำแหน่ง และขอบประด้านตรงข้ามใช้นับจำนวน module
ความยากในการอ่านขึ้นกับจำนวนพิกเซลที่ตกบน module หนึ่งช่อง
ซึ่งนิยามได้เป็น
\begin{equation}
\mathrm{ppm} = \frac{w_{\mathrm{code}}}{n_{\mathrm{mod}}}\, r
\label{eq:ppm}
\end{equation}
เมื่อ $w_{\mathrm{code}}$ คือความกว้างของรหัส (มม.)
$n_{\mathrm{mod}}$ คือจำนวน module ต่อด้าน
และ $r$ คือจำนวนพิกเซลต่อมิลลิเมตรบนระนาบฉลาก

\textbf{แบบจำลองกล้องและ homography}
จุดบนระนาบฉลาก $\mathbf{x} = (x, y, 1)^\top$ ในหน่วยมิลลิเมตร
ฉายไปยังพิกเซล $\mathbf{u}$ ในภาพผ่าน homography $\mathbf{H}$ ขนาด $3\times3$
\begin{equation}
\lambda\,\mathbf{u} = \mathbf{H}\,\mathbf{x}
\label{eq:homography}
\end{equation}
ความสัมพันธ์นี้ใช้ทั้งในการสร้างภาพสังเคราะห์ที่รู้ตำแหน่ง 4~มุมแน่นอน
และในการดัดภาพรหัสที่เอียงให้กลับเป็นสี่เหลี่ยมจัตุรัส
โดยหา $\mathbf{H}$ จากมุม 4~จุดของรหัส

\textbf{Instance segmentation}
เป็นงานที่ทำนายทั้งตำแหน่งและ mask ระดับพิกเซลของวัตถุแต่ละชิ้น
Mask R-CNN~\cite{he2017maskrcnn} เป็นแบบ two-stage
ที่เพิ่มส่วนทำนาย mask ต่อจาก Faster R-CNN
YOLO11~\cite{khanam2024yolov11} และ YOLO26~\cite{sapkota2025yolo26}
เป็นแบบ one-stage ที่เน้นความเร็ว
และ RF-DETR~\cite{robinson2026rfdetr} เป็นแบบ transformer
ที่ใช้ neural architecture search หาสถาปัตยกรรมที่สมดุลระหว่างความแม่นยำและความเร็ว

\textbf{Domain randomization}
การสุ่มพารามิเตอร์ของตัวสร้างภาพ เช่น แสง พื้นหลัง และความเสื่อมของภาพ
ช่วยให้โมเดลที่ฝึกด้วยภาพสังเคราะห์นำไปใช้กับภาพจริงได้ดีขึ้น~\cite{tremblay2018domain}

\textbf{Knowledge distillation}
เป็นการฝึกโมเดลขนาดเล็ก (student) ให้เลียนแบบผลลัพธ์ของโมเดลขนาดใหญ่ (teacher)~\cite{hinton2015distilling}
เพื่อให้ได้โมเดลที่เร็วขึ้นโดยเสียความแม่นยำน้อย
รูปแบบทั่วไปของฟังก์ชันสูญเสียคือ
\begin{equation}
\mathcal{L} = (1-\alpha)\,\mathcal{L}_{\mathrm{task}} + \alpha\,\mathcal{L}_{\mathrm{KD}}
\label{eq:kd}
\end{equation}
เมื่อ $\mathcal{L}_{\mathrm{task}}$ คือ loss จากป้ายกำกับ
$\mathcal{L}_{\mathrm{KD}}$ คือ loss ที่วัดความต่างระหว่างผลของ student กับ teacher
และ $\alpha$ คือน้ำหนัก

\subsection{การสังเคราะห์งานวิจัย (Literature Synthesis)}

งานที่เกี่ยวข้องปี 2019--2026 จำนวน 9~งาน
แบ่งได้เป็น 3~กลุ่มตามวิธีบอกตำแหน่งรหัส
กลุ่มที่~1 บอกตำแหน่งด้วยกรอบ~\cite{ershova2024yolobarcode,pandya2025yolov8barcode,alsulami2026uav}
กลุ่มที่~2 บอกตำแหน่งด้วย mask~\cite{zharkov2019universal,quenum2021uhr,yuan2019murcnn}
และกลุ่มที่~3 หามุมหรือขอบแล้วจัดรหัสให้ตรงก่อนอ่าน~\cite{rybakova2025datamatrix,liu2024lshaped,corpuz2025qr}
รายละเอียดแสดงในตารางที่~\ref{tab:synthesis_matrix}
ตัวเลขในตารางมาจากชุดข้อมูลและตัววัดของแต่ละงาน จึงไม่สามารถเปรียบเทียบข้ามงานโดยตรงได้

\begin{table}[H]
\centering
\caption{ตารางสังเคราะห์และเปรียบเทียบงานวิจัย}
\label{tab:synthesis_matrix}

\begin{adjustbox}{width=\textwidth}
\footnotesize

\begin{tabular}{|
>{\centering\arraybackslash}p{1.0cm}|
p{2.5cm}|
p{2.5cm}|
p{2.3cm}|
p{2.5cm}|
p{3.0cm}|}
\hline

\multicolumn{1}{|c|}{\textbf{งานวิจัย}} &
\multicolumn{1}{c|}{\textbf{วิธีการ}} &
\multicolumn{1}{c|}{\textbf{ชุดข้อมูล}} &
\multicolumn{1}{c|}{\textbf{ตัววัด}} &
\multicolumn{1}{c|}{\textbf{ผลลัพธ์}} &
\multicolumn{1}{c|}{\textbf{ข้อจำกัด}}
\\
\hline

\cite{ershova2024yolobarcode} &
YOLO ตรวจจับและบอกชนิดรหัส 1D/2D หลายชนิด รวม DataMatrix &
ZVZ-real, ZVZ-synth และ SE-barcode 1,813~ภาพ มีหลายรหัสต่อภาพ &
F1-score ของการตรวจจับ, เวลา &
F1 98.6\% บน ZVZ-real (งานเดิม 88.9\%) เร็วขึ้น 4~เท่า &
ทำนายแค่กรอบไม่เอียง ไม่ได้ 4~มุม และไม่วัดการอ่านรหัส
\\
\hline

\cite{pandya2025yolov8barcode} &
เทียบ YOLOv8 ขนาด n, s, m &
Kaggle 31,078~ภาพ (1D และ QR) &
Precision, Recall, F1, Accuracy &
n 88.95\%, s 97.10\%, m 94.10\% &
วัดแค่การตรวจจับ ไม่มี DataMatrix รายละเอียดการทดลองน้อย
\\
\hline

\cite{alsulami2026uav} &
YOLOv8 หากรอบ แล้วถอดรหัสด้วย OpenCV &
31,510~ภาพ (1D และ QR) และภาพจำลองความเสื่อม 300~ภาพ &
mAP@0.5, อัตราการถอดรหัสแยกตามสภาพภาพ &
mAP@0.5 0.922 $\pm$ 0.0026 (6~รอบ) อ่านได้ 98.5\% (แสงดี) ถึง 72.1\% (เบลอ) &
ถอดรหัสจากกรอบสี่เหลี่ยม ไม่แยกผลตามความละเอียด ไม่มี DataMatrix
\\
\hline

\cite{zharkov2019universal} &
semantic segmentation ขนาดเล็ก แล้วแปลง mask เป็นกรอบเอียง &
17,000~ภาพ 21~ชนิด รวม DataMatrix &
detection rate, เวลา &
detection rate 0.989 บน ArTe-Lab &
ได้กรอบเอียง ไม่วัดการอ่าน
\\
\hline

\cite{quenum2021uhr} &
RPN เลือกบริเวณ แล้ว Y-Net ทำ mask และฟิตกรอบ &
ภาพสังเคราะห์ใหญ่เกิน 10k$\times$10k หลายชนิด รวม Data Matrix &
mAP, AR, เวลา &
mAP 0.937 (YOLOv4 0.882, Mask R-CNN 0.466) &
จบที่กรอบ ไม่วัดการอ่าน
\\
\hline

\cite{yuan2019murcnn} &
Mask R-CNN เพิ่มกิ่ง UNet และ dice loss &
ภาพโรงงาน 14,000~ภาพ 18,409~QR และชุดอื่น &
mask AP, box AP &
mask AP 17.38 (Mask R-CNN 8.66) &
มีแต่ QR และไม่มีการทดลองอ่าน
\\
\hline

\cite{rybakova2025datamatrix} &
Fast Hough transform หาขอบรูปตัว L ได้ 4~มุม แล้วแบ่งช่อง module ซ้ำ &
SE-DMTX-SYN-1000 (กึ่งสังเคราะห์ 1,000~รหัส) &
mIoU ของตำแหน่ง, อัตราการอ่าน, เวลาบน ARM &
mIoU 0.889 อ่านได้ 90.3\% (zxing-cpp 85.8\%) ใช้ 35~ms &
รหัสเดียวต่อภาพและต้องกินพื้นที่ $\geq$ 25\% ของภาพ สมมติว่ามีตัวหา ROI ให้ก่อน
\\
\hline

\cite{liu2024lshaped} &
YOLOv5s หาตำแหน่งคร่าว หาขอบ L เส้นประ แล้วอ่านด้วย libdmtx &
ทดสอบ 265~ภาพ 5~แบบ DataMatrix &
recognition accuracy rate, เวลา &
วิธีเต็มอ่านได้ 90.2\% รหัสหมุน 85\% &
ไม่ระบุจำนวนรหัสต่อภาพ ไม่แยกตามความละเอียด
\\
\hline

\cite{corpuz2025qr} &
YOLOv5 หากรอบ Canny หา 4~มุม แล้วแก้มุมมองก่อนอ่าน &
Roboflow 1,338~ภาพ QR &
จำนวนที่อ่านได้, มุมสูงสุด, เวลา &
อ่านได้ 459~รหัส (ZBar 352, OpenCV 214) &
ผู้เขียนระบุว่ายังต้องรองรับหลาย QR ในภาพเดียว
\\
\hline

\end{tabular}

\end{adjustbox}
\end{table}

ข้อมูลในตารางตรวจสอบกับฉบับเต็มของทุกงานแล้ว

%-----------------------------------------------%
% 4. Research Gap and Questions
%-----------------------------------------------%

\section{ช่องว่างและคำถามวิจัย
(Research Gap and Research Questions)}

\subsection{ช่องว่างของงานวิจัย (Research Gap)}

จากตารางที่~\ref{tab:synthesis_matrix} แต่ละกลุ่มแก้ปัญหาได้ทีละส่วน
กลุ่มที่~1 ตรวจจับได้แม่นแต่กรอบไม่แก้ความเอียง
กลุ่มที่~2 ได้รูปทรงจริงของรหัสแต่แปลงกลับเป็นกรอบและไม่วัดการอ่าน
ส่วนกลุ่มที่~3 จัดรหัสให้ตรงแล้วอ่านได้ดีขึ้น
แต่ยังไม่ได้ทดสอบรหัสเล็กหลายรหัสในภาพเดียว

\textbf{ช่องว่างร่วม (Common Gap)} คือ
ในงานที่ทบทวน ยังไม่มีงานที่ใช้ mask หา 4~มุมของรหัส DataMatrix ขนาดเล็กหลายรหัสในภาพความละเอียดสูงภาพเดียว
แล้วประเมินผลถึงขั้นอ่านได้จริงและอ่านผิด โดยแยกตาม px/module
ซึ่งเป็นค่าที่บอกว่าภาพต้องละเอียดเท่าใด
ช่องว่างนี้เป็นทั้ง Evaluation Gap, Dataset Gap และ Application Gap

เหตุผลที่ยังไม่มีงานลักษณะนี้มี 3~ข้อ
ข้อแรก การติดป้าย 4~มุมของรหัสขนาดเล็กด้วยมือทำได้ยากและไม่แม่นยำ
ข้อที่สอง ภาพฉลากสไลด์จริงมีข้อมูลผู้ป่วย จึงเผยแพร่เป็นชุดข้อมูลไม่ได้
และข้อที่สาม งานบาร์โค้ดส่วนใหญ่มาจากโลจิสติกส์และค้าปลีก ซึ่งรหัสมีขนาดใหญ่และอ่านทีละอันได้

\subsection{คำถามวิจัย (Research Questions)}

\begin{description}
\item[\textbf{RQ1}:] ใช้โมเดล instance segmentation ใดถึงจะเหมาะสม
(Mask R-CNN, YOLO11n-seg, YOLO26n-seg และ RF-DETR-Seg ฝึกด้วยเวลาเท่ากัน)
\item[\textbf{RQ2} (คำถามหลัก):] เตรียมภาพรหัสแบบใดถึงจะเหมาะสม
(ป้อนทั้งภาพ ตัดตามกรอบ ดัดจาก 4~มุม หรือดัดแล้วใช้การตัดตามกรอบเป็นขั้นสำรอง)
\item[\textbf{RQ3}:] ความละเอียดของภาพเท่าใด (px/module) ถึงจะเหมาะสม
\end{description}

\textbf{สมมติฐานหลัก (RQ2)}
$H_0$: การดัดภาพจาก 4~มุมที่ทำนาย โดยใช้การตัดกรอบเป็นขั้นสำรอง
ไม่ทำให้อัตราการอ่านสูงกว่าการตัดตามกรอบสี่เหลี่ยมเพียงอย่างเดียวอย่างมีนัยสำคัญทางสถิติ
$H_1$: การดัดภาพดังกล่าวทำให้อัตราการอ่านสูงกว่าอย่างมีนัยสำคัญทางสถิติ

%-----------------------------------------------%
% 5. Research Objectives
%-----------------------------------------------%

\section{วัตถุประสงค์ของโครงงาน (Research Objectives)}

\begin{enumerate}
    \item เพื่อสร้างชุดข้อมูลสังเคราะห์ของฉลากสไลด์พยาธิวิทยาที่มีรหัส DataMatrix หลายรหัสต่อภาพ
    พร้อมตำแหน่ง 4~มุมที่แม่นยำ และควบคุมความละเอียดได้ในช่วง 2--8~px/module

    \item เพื่อหาโมเดลที่เหมาะสมที่สุดสำหรับหารหัสเล็กหลายรหัสในภาพเดียว ให้อ่านได้มากและเร็วพอใช้งาน

    \item เพื่อหาวิธีเตรียมภาพที่อ่านรหัสได้มากที่สุด โดยไม่เพิ่มการอ่านผิด

    \item เพื่อหาความละเอียดที่เหมาะสมต่อการอ่านรหัส
    และพัฒนาระบบต้นแบบที่ประกอบด้วยแท่นถ่ายภาพจริง โมเดลขนาดเล็กจาก Knowledge Distillation และ Web app
\end{enumerate}

%-----------------------------------------------%
% 6. Methodology
%-----------------------------------------------%

\section{วิธีการดำเนินงานวิจัย (Methodology)}

\subsection{ภาพรวมวิธีการที่เสนอ
(Proposed Method Overview)}

ระบบรับภาพสไลด์ทั้งแถวจากแท่นถ่ายภาพ (Input)
แล้วทำงาน 5~ขั้นตอน ดังนี้
(1)~โมเดล instance segmentation ทำนาย mask ของรหัสแต่ละอัน
(2)~หาเส้นขอบของ mask แล้วประมาณเป็นรูปสี่เหลี่ยม 4~มุม
(3)~ดัดภาพรหัสจาก 4~มุมให้เป็นสี่เหลี่ยมจัตุรัสด้วย homography ตามสมการที่~\ref{eq:homography}
(4)~ส่งภาพที่ดัดแล้วให้ zxing-cpp อ่าน
และ (5)~ถ้าอ่านไม่ได้ ตัดภาพตามกรอบสี่เหลี่ยมของ mask (ขยายขอบ 25\%) แล้วอ่านอีกครั้ง
ผลลัพธ์ (Output) คือข้อความในรหัสของสไลด์ทุกแผ่นในภาพ
ตัวอย่างแสดงในภาพที่~\ref{fig:pipeline}

\begin{figure}[H]
\centering
\includegraphics[width=\linewidth]{figures/strip_raw.png}\\[2pt]
\includegraphics[width=\linewidth]{figures/strip_mask.png}\\[2pt]
\includegraphics[width=\linewidth]{figures/strip_read.png}
\caption{ตัวอย่างการทำงานบนภาพสังเคราะห์: ภาพนำเข้า (บน) mask ที่ YOLO11n-seg ทำนาย (กลาง) และผลการอ่าน (ล่าง) สีเขียวคืออ่านได้จากการดัดภาพจาก 4~มุม}
\label{fig:pipeline}
\end{figure}

\subsection{รายละเอียดวิธีการที่เสนอ
(Proposed Method)}

\textbf{การหา 4~มุมจาก mask}
หาเส้นขอบที่ใหญ่ที่สุดของ mask แล้วประมาณเป็นรูปหลายเหลี่ยมด้วย Douglas--Peucker
โดยค่อยเพิ่มค่าความคลาดเคลื่อนที่ยอมได้จนได้รูปสี่เหลี่ยมนูน
ถ้าไม่ได้ ใช้สี่เหลี่ยมพื้นที่น้อยที่สุดที่ครอบ mask แทน
ความคลาดเคลื่อนของมุมวัดเทียบกับมุมจริง $\mathbf{g}_i$
โดยเลือกการเรียงมุมที่ดีที่สุด เพราะ mask ไม่บอกทิศของรหัส
\begin{equation}
e = \min_{k \in \{0,1,2,3\}} \frac{1}{4}\sum_{i=0}^{3}
\left\lVert \mathbf{q}_{(i+k) \bmod 4} - \mathbf{g}_i \right\rVert_2
\label{eq:corner}
\end{equation}

\textbf{การดัดภาพและวิธีผสม}
คำนวณ homography จาก 4~มุมที่ทำนายไปยังสี่เหลี่ยมจัตุรัสขนาดคงที่ (8~พิกเซลต่อ module)
แล้วส่งภาพที่ดัดแล้วให้ตัวอ่าน
วิธีผสมใช้ผลจากการดัดภาพก่อน และใช้การตัดตามกรอบเฉพาะเมื่อการดัดภาพอ่านไม่ได้
\begin{equation}
t_{\mathrm{hyb}} =
\begin{cases}
t_{\mathrm{warp}} & \text{ถ้าการดัดภาพอ่านได้} \\
t_{\mathrm{crop}} & \text{ถ้าการดัดภาพอ่านไม่ได้}
\end{cases}
\label{eq:hybrid}
\end{equation}

\textbf{ตัวอ่าน}
ใช้ zxing-cpp~\cite{zxingcpp} ตัวเดียวกันทุกวิธี
โดยลอง binarizer แบบ LocalAverage ก่อน และใช้ GlobalHistogram เมื่ออ่านไม่ได้
ตัวอ่านจึงเป็นตัวแปรควบคุม ไม่ใช่สิ่งที่ศึกษา

\textbf{Knowledge Distillation (ภาคการศึกษาถัดไป)}
ฝึกโมเดลขนาดเล็กให้เลียนแบบโมเดลที่แม่นกว่าตามสมการที่~\ref{eq:kd}
เพื่อให้รันบนคอมพิวเตอร์ในห้องปฏิบัติการได้เร็ว
คู่ teacher กับ student และรูปแบบของ $\mathcal{L}_{\mathrm{KD}}$ จะกำหนดในภาคการศึกษาถัดไป

\textbf{Web app และแท่นถ่ายภาพ (ภาคการศึกษาถัดไป)}
แท่นถ่ายภาพมีฐาน PLA ขนาด 380$\times$150~มม.
มีกรอบวางสไลด์ 10~ช่อง และขาจับกล้องที่ยึดระยะคงที่ (ภาพที่~\ref{fig:rig})
Web app จะรับภาพจากกล้องหรือการอัปโหลด แล้วแสดงผลการอ่านของทุกแผ่น
พร้อมบอกแผ่นที่อ่านไม่ได้ให้เจ้าหน้าที่สแกนเอง

\begin{figure}[H]
\centering
\includegraphics[width=0.45\linewidth]{figures/rig_iso_crop.png}
\caption{แบบ CAD ของแท่นถ่ายภาพที่จะสร้างในภาคการศึกษาถัดไป}
\label{fig:rig}
\end{figure}

\subsection{ข้อมูลและการเตรียมข้อมูล
(Datasets and Data Preparation)}

ใช้ภาพสังเคราะห์ที่สร้างขึ้นเองทั้งหมด
ตัวสร้างภาพวางฉลากในหน่วยมิลลิเมตรแล้วฉายผ่านแบบจำลองกล้องตามสมการที่~\ref{eq:homography}
จึงได้ตำแหน่ง 4~มุมของทุกรหัสแบบแม่นยำโดยไม่ต้องติดป้ายด้วยมือ
ข้อความบนฉลากสุ่มขึ้นใหม่ทั้งหมด ไม่ได้มาจากภาพจริง
ภาพมีขนาด 3264$\times$2448~พิกเซล
มีสไลด์ 1--10~แผ่นต่อภาพ และสุ่มความละเอียด 2--8~px/module

\begin{itemize}
    \item \textbf{ขนาดข้อมูล:} 3,600~ภาพ 19,506~รหัส
    แบ่งเป็นภาพจากแท่น 3,000~ภาพ (train 2,100 / val 450 / test 450)
    และภาพถือกล้องด้วยมือ 600~ภาพ ซึ่งใช้ทดสอบอย่างเดียว

    \item \textbf{Data Augmentation:} สุ่มความเสื่อมของภาพตามแนวคิด domain randomization~\cite{tremblay2018domain}
    ได้แก่ คราบบนฉลาก ฉลากยับ อุณหภูมิสี แสงไม่สม่ำเสมอ เงา แสงสะท้อน
    ภาพไม่ชัด ภาพสั่น vignetting noise และการบีบอัด JPEG

    \item \textbf{Data Splitting Strategy และการป้องกัน Data Leakage:}
    แบ่งข้อมูลเป็นรายภาพ ไม่ใช่รายรหัส รหัสในภาพเดียวกันจึงอยู่ชุดเดียวกันเสมอ
    แบ่งให้ทุกชุดมีสัดส่วนความละเอียดเท่ากัน
    และใช้ชุด test เฉพาะการประเมินผลขั้นสุดท้าย

    \item \textbf{รหัสที่หลุดขอบภาพ:} รหัสที่เห็นน้อยกว่าครึ่งถูกตัดออกจากการให้คะแนน
\end{itemize}

ภาคการศึกษาถัดไปจะเก็บภาพจริงจากแท่นถ่ายภาพในห้องปฏิบัติการเพื่อใช้ทดสอบ
ภาพจริงจะเก็บไว้ในเครื่องเท่านั้น ไม่เผยแพร่ และใช้เท่าที่จำเป็นตามหลัก data minimization

\subsection{โมเดลเปรียบเทียบ
(Baselines)}

\begin{itemize}
    \item \textbf{สถาปัตยกรรม:} Mask R-CNN R50-FPN~\cite{he2017maskrcnn} (two-stage),
    YOLO11n-seg~\cite{khanam2024yolov11} และ YOLO26n-seg~\cite{sapkota2025yolo26} (one-stage)
    และ RF-DETR-Seg Nano~\cite{robinson2026rfdetr} (transformer, SOTA)

    \item \textbf{วิธีเตรียมภาพที่เป็น baseline:} ป้อนภาพทั้งภาพให้ zxing-cpp โดยตรง
    และการตัดตามกรอบสี่เหลี่ยม~\cite{alsulami2026uav} ซึ่งรันเองบนข้อมูลชุดเดียวกัน

    \item \textbf{เพดานบน (oracle):} ใช้ 4~มุมจริงแทนมุมที่ทำนาย
    แล้วผ่านขั้นตอนเดียวกันทุกอย่าง เพื่อแยกความผิดพลาดของโมเดลออกจากข้อจำกัดของตัวอ่าน
\end{itemize}

\subsection{สภาพแวดล้อมและการประเมินผล
(Experimental Settings and Evaluation)}

\textbf{Hardware และ Software}
GPU NVIDIA RTX~5070 (VRAM 12~GB) RAM 32~GB
ฝึกบน Ubuntu 24.04 ใน WSL2 ด้วย Python และ PyTorch
ใช้ Ultralytics สำหรับ YOLO, torchvision สำหรับ Mask R-CNN
และ rfdetr สำหรับ RF-DETR

\textbf{Hyperparameters และ Random Seed}
ทุกโมเดลใช้ภาพนำเข้าด้านยาวประมาณ 1,280~พิกเซล (RF-DETR ใช้ 1,272~พิกเซล เพราะต้องหารด้วย 12 ลงตัว)
เวลาฝึก 60~นาทีเท่ากัน (รวมการ validate)
และ random seed เดียวกัน (seed~0)
Mask R-CNN ใช้ batch~2 learning rate 0.01
YOLO ใช้ batch~8
และ RF-DETR ใช้ batch~1 สะสม gradient ให้เทียบเท่า batch~8 learning rate 0.0001
ภาคการศึกษาถัดไปจะฝึกซ้ำ 3~seed

\textbf{Metrics}
ตัวชี้วัดหลักคืออัตราการอ่านได้ถูกต้อง
\begin{equation}
R = \frac{1}{N}\sum_{j=1}^{N} \mathbf{1}\!\left[\hat{t}_j = t_j\right]
\label{eq:readrate}
\end{equation}
เมื่อ $t_j$ คือข้อความจริงของรหัสที่ $j$ และ $\hat{t}_j$ คือข้อความที่อ่านได้
นอกจากนี้วัดจำนวนครั้งที่อ่านผิด
คือข้อความที่อ่านได้ไม่ตรงกับรหัสที่อยู่ใต้ mask
(รหัสที่ทับกับ mask อย่างน้อย 50\% ของพื้นที่ที่เล็กกว่า)
รวมถึงความคลาดเคลื่อนของมุมตามสมการที่~\ref{eq:corner}
AP50 และ AP50:95 ของ mask และเวลาประมวลผลต่อภาพบน GPU และ CPU
ทุกตัวชี้วัดรายงานแยกตามช่วง px/module
โดยเน้นช่วง 3.5--4.0 (กล้อง 8~ล้านพิกเซล) และ 4.5--5.0 (กล้อง 12~ล้านพิกเซล)

\textbf{วิธีการทางสถิติ}
เนื่องจากฝึกแต่ละโมเดล 1~ครั้ง จึงใช้การทดสอบแบบจับคู่บนรหัสชุดเดียวกัน
ช่วงความเชื่อมั่น 95\% คำนวณด้วย cluster bootstrap 2,000~รอบ โดยสุ่มทั้งภาพ
เพราะรหัสในภาพเดียวกันได้แสงและความเสื่อมเดียวกัน จึงไม่เป็นอิสระต่อกัน
เปรียบเทียบสองวิธีด้วย exact McNemar test
\begin{equation}
p = \min\!\left(1,\; 2\sum_{i=0}^{\min(b,c)} \binom{b+c}{i} 2^{-(b+c)}\right)
\label{eq:mcnemar}
\end{equation}
เมื่อ $b$ และ $c$ คือจำนวนรหัสที่อ่านได้เฉพาะวิธีแรกและเฉพาะวิธีที่สอง
ปรับค่า $p$ ด้วยวิธี Holm เมื่อเทียบหลายคู่
และใช้ Cochran's Q เมื่อเทียบ 4~โมเดลพร้อมกัน ที่ระดับนัยสำคัญ $\alpha = 0.05$

\textbf{ผลการทดลองเบื้องต้น (ภาพสังเคราะห์)}
ภาพที่~\ref{fig:readrate} แสดงอัตราการอ่านของ YOLO11n-seg แยกตามความละเอียด
ที่ช่วงกล้อง 8~ล้านพิกเซล การป้อนทั้งภาพอ่านได้ 64.1\%
การตัดตามกรอบได้ 88.5\% และวิธีผสมได้ 95.7\%
ที่ช่วงกล้อง 12~ล้านพิกเซล วิธีผสมได้ 99.4\%
และไม่มีการอ่านผิดเลยจากการอ่านได้ทั้งหมด 4,754~ครั้ง
เมื่อรวมทุกช่วง วิธีผสมสูงกว่าการตัดตามกรอบ 4.0--6.3~จุดเปอร์เซ็นต์ในทุกโมเดล ($p < 0.001$)
แต่ที่ช่วง 12~ล้านพิกเซลยังไม่มีนัยสำคัญ
ส่วน YOLO11n-seg, YOLO26n-seg และ Mask R-CNN อ่านได้ไม่ต่างกันทางสถิติ
โดย RF-DETR ต่ำกว่า 1.4--2.2~จุด เพราะใน 60~นาทีฝึกได้เพียง 2~epoch
ข้อควรระวังคือวิธีผสมใช้การตัดกรอบเป็นขั้นสำรอง จึงไม่มีทางอ่านได้น้อยกว่า baseline
จึงรายงานขนาดของผลต่างพร้อมช่วงความเชื่อมั่นเป็นหลัก

\begin{figure}[H]
\centering
\includegraphics[width=\linewidth]{figures/read_rate_yolo11n_rig.pdf}
\caption{อัตราการอ่านรหัสตามความละเอียดของ YOLO11n-seg บนชุดทดสอบจากแท่นถ่ายภาพ (ภาพสังเคราะห์ 450~ภาพ) แถบสีเทาคือช่วงความละเอียดของกล้อง 8 และ 12~ล้านพิกเซล}
\label{fig:readrate}
\end{figure}

%-----------------------------------------------%
% 7. Scope and Limitations
%-----------------------------------------------%

\section{ขอบเขตและข้อจำกัดของการวิจัย
(Scope and Limitations)}

\begin{enumerate}
    \item \textbf{ขอบเขต:} ศึกษารหัส DataMatrix 18$\times$18~module ขนาด 6.5~มม.
    บนฉลากสไลด์สูงสุด 10~แผ่นต่อภาพ
    เปรียบเทียบโมเดล 4~สถาปัตยกรรมและวิธีเตรียมภาพ 4~แบบ
    ใช้ภาพสังเคราะห์ในการฝึกและทดสอบ
    และพัฒนาระบบต้นแบบที่ประกอบด้วยแท่นถ่ายภาพ โมเดลจาก Knowledge Distillation และ Web app
    \item \textbf{ข้อจำกัด:} ผลทั้งหมดในเอกสารนี้มาจากภาพสังเคราะห์และยังไม่ได้ทดสอบกับภาพจริง
    ขนาดรหัสและระยะห่างของสไลด์เป็นค่าที่ตั้งไว้ ยังต้องยืนยันกับฉลากจริง
    แต่ละโมเดลฝึกเพียง 1~seed จึงยังไม่รวมความแปรปรวนจากการฝึก
    ใช้ GPU เพียงเครื่องเดียว
    และ YOLO มี license แบบ AGPL-3.0 ซึ่งมีเงื่อนไขเมื่อนำไปใช้งานจริง
\end{enumerate}

%-----------------------------------------------%
% 8. Expected Results
%-----------------------------------------------%

\section{ผลลัพธ์ที่คาดว่าจะได้รับ (Expected Results)}

คาดว่าจะได้รับผลลัพธ์จากการดำเนินโครงงาน ดังนี้

\begin{itemize}

    \item \textbf{โมเดลหรือวิธีการที่พัฒนา:}
    ได้ระบบหาตำแหน่ง 4~มุมและอ่านรหัส DataMatrix หลายรหัสจากภาพเดียว
    และโมเดลขนาดเล็กจาก Knowledge Distillation ที่รันบนคอมพิวเตอร์ในห้องปฏิบัติการได้

    \item \textbf{ผลการประเมินประสิทธิภาพ:}
    ได้ผลเปรียบเทียบอัตราการอ่าน การอ่านผิด และเวลาประมวลผลของโมเดล 4~สถาปัตยกรรม
    รวม SOTA อย่าง RF-DETR แยกตามความละเอียด ทั้งบนภาพสังเคราะห์และภาพจริง

    \item \textbf{ผลการวิเคราะห์โมเดล:}
    ได้ผล Ablation Study ของวิธีเตรียมภาพแต่ละแบบ
    ซึ่งแสดงว่าการดัดภาพจาก 4~มุมช่วยเพิ่มการอ่านได้เท่าใดและที่ความละเอียดใด

    \item \textbf{องค์ความรู้และข้อค้นพบ:}
    ได้ความละเอียดขั้นต่ำที่อ่านรหัสได้เพียงพอ และข้อแนะนำในการถ่ายภาพ

    \item \textbf{ซอฟต์แวร์และเอกสารประกอบ:}
    ได้ Source Code ของตัวสร้างภาพสังเคราะห์ การฝึก และการประเมินผล
    แบบ CAD ของแท่นถ่ายภาพ และ Web app ต้นแบบ
    เพื่อสนับสนุนการทำซ้ำผลการทดลอง (Reproducibility)

    \item \textbf{ผลงานทางวิชาการ:}
    ได้รายงานผลการวิจัยที่สามารถนำไปเผยแพร่ในงานประชุมวิชาการที่เหมาะสม

\end{itemize}

%-----------------------------------------------%
% 9. Research Location
%-----------------------------------------------%

\section{สถานที่ทำวิจัย (Research Location)}

วิทยาลัยการคอมพิวเตอร์ มหาวิทยาลัยขอนแก่น
และห้องปฏิบัติการพยาธิวิทยาที่ร่วมโครงการ (สำหรับเก็บภาพจริง)

%-----------------------------------------------%
% 10. Timeline
%-----------------------------------------------%

\section{แผนและระยะเวลาดำเนินการ
(Timeline / Gantt Chart)}

แสดงแผนการดำเนินงานและระยะเวลาของแต่ละกิจกรรม
ตลอดระยะเวลาของโครงงาน

\textbf{แผนการดำเนินงานวิจัย}
แผนการดำเนินงานประกอบด้วยขั้นตอนและระยะเวลา
ดังแสดงในตารางที่ \ref{tab:timeline}

\begin{table}[H]
\centering
\caption{แผนการดำเนินงานวิจัย}
\label{tab:timeline}

\begin{adjustbox}{width=\textwidth}
\begin{tabular}{|p{0.4\textwidth}*{12}{|>{\centering\arraybackslash}p{0.03\textwidth}}|}
\hline

\multirow{2}{*}{
\begin{minipage}[b]{0.38\textwidth}
\centering
\textbf{การดำเนินงาน}
\end{minipage}
}
&
\multicolumn{6}{c|}{ปีการศึกษา 2569}
&
\multicolumn{6}{c|}{ปีการศึกษา 2570}
\\
\cline{2-13}

& มิ.ย. & ก.ค. & ส.ค. & ก.ย. & ต.ค. & พ.ย.
& ธ.ค. & ม.ค. & ก.พ. & มี.ค. & เม.ย. & พ.ค.
\\
\hline

1. กำหนดหัวข้อวิจัย
\on[1] \off[11]
\\ \hline

2. ทบทวนวรรณกรรมที่เกี่ยวข้อง
\on[2] \off[10]
\\ \hline

3. กำหนดคำถามวิจัยและวัตถุประสงค์
\off[1] \on[1] \off[10]
\\ \hline

4. ออกแบบวิธีการทดลอง
\off[2] \on[1] \off[9]
\\ \hline

5. รวบรวมและเตรียมข้อมูล
\off[2] \on[2] \off[8]
\\ \hline

6. พัฒนาและฝึกโมเดล
\off[4] \on[5] \off[3]
\\ \hline

7. ทดสอบโมเดลและประเมินผล
\off[8] \on[2] \off[2]
\\ \hline

8. เปรียบเทียบประสิทธิภาพ
\off[9] \on[2] \off[1]
\\ \hline

9. วิเคราะห์และสรุปผลการทดลอง
\off[10] \on[1] \off[1]
\\ \hline

10. นำเสนอโครงงาน
\off[11] \on[1]
\\ \hline

\end{tabular}
\end{adjustbox}
\end{table}


%-----------------------------------------------%
% 11. Resources and Budget
%-----------------------------------------------%

\section{ทรัพยากรและงบประมาณ
(Resources and Budget)}

\subsection{ทรัพยากรโครงงาน (Project Resources)}

\begin{itemize}

    \item คอมพิวเตอร์ที่มี GPU NVIDIA RTX~5070 (VRAM 12~GB) และ RAM 32~GB
    สำหรับสร้างข้อมูล ฝึก และทดสอบโมเดล

    \item กล้องความละเอียด 8--12~ล้านพิกเซล
    และแท่นถ่ายภาพที่พิมพ์ด้วยเครื่องพิมพ์ 3~มิติ (PLA) ร่วมกับโครงอะลูมิเนียม

    \item ภาษา Python บน Ubuntu 24.04 (WSL2) และ Visual Studio Code

    \item Library และ Framework ได้แก่ PyTorch, torchvision, Ultralytics, rfdetr,
    OpenCV, NumPy, pycocotools, zxing-cpp และ CadQuery สำหรับงาน CAD

    \item ชุดข้อมูลสังเคราะห์ที่สร้างขึ้นเอง และภาพจริงจากห้องปฏิบัติการพยาธิวิทยา (ภาคการศึกษาถัดไป)

\end{itemize}

\subsection{ค่าใช้จ่ายที่จำเป็น (Budget) (ถ้ามี)}

ค่าใช้จ่ายต่อไปนี้เป็นการประมาณการ ยังไม่ได้ตรวจราคาจากผู้ขายจริง
ส่วนคอมพิวเตอร์และ GPU ใช้เครื่องที่มีอยู่แล้ว จึงไม่มีค่าใช้จ่ายเพิ่ม
รายละเอียดแสดงในตารางที่~\ref{tab:budget}

\begin{table}[H]
\centering
\caption{ประมาณการค่าใช้จ่าย}
\label{tab:budget}
\begin{tabular}{|p{0.62\textwidth}|>{\raggedleft\arraybackslash}p{0.2\textwidth}|}
\hline
\multicolumn{1}{|c|}{\textbf{รายการ}} & \multicolumn{1}{c|}{\textbf{ประมาณการ (บาท)}} \\ \hline
กล้อง USB ความละเอียด 12~ล้านพิกเซลพร้อมเลนส์ & 2,500 \\ \hline
เส้นพลาสติก PLA สำหรับพิมพ์ 3~มิติ 1~กก. & 600 \\ \hline
อะลูมิเนียมโปรไฟล์ 2020 ยาว 400~มม. น็อต สกรู และอุปกรณ์ยึดกล้อง & 500 \\ \hline
ค่าใช้จ่ายสำรอง (10\%) & 360 \\ \hline
\textbf{รวม} & \textbf{3,960} \\ \hline
\end{tabular}
\end{table}

%-----------------------------------------------%
% 12. References
%-----------------------------------------------%

\renewcommand{\refname}{บรรณานุกรม (References)}
\bibliographystyle{unsrt}
\bibliography{reference}

%-----------------------------------------------%
% ลายเซ็น
%-----------------------------------------------%

\begin{flushright}

\begin{tabular}{>{\centering\arraybackslash}p{0.4\textwidth}}

\vspace{0.5em}

ลงชื่อผู้ทำโครงงาน \dotfill
\\
(\firstStudentPrefixTH\firstStudentNameTH)
\\
วันที่ \dotfill/\dotfill/\dotfill
\\[1em]

ลงชื่อผู้ทำโครงงาน \dotfill
\\
(\secondStudentPrefixTH\secondStudentNameTH)
\\
วันที่ \dotfill/\dotfill/\dotfill

\end{tabular}

\end{flushright}



\end{document}
